\uuid{OVpA}
\titre{Changements de variables}
\auteur{hanani ; mzari}
\organisation{exo7}
\licence{CC BY-NC-SA 4.0}
\datecreate{2020-01-20}
\niveau{L1}
\module{Analyse}
\chapitre{Calcul d'intégrales}
\sousChapitre{Autre}
\difficulte{3}
\theme{Calculs d'intégrales, changements de variables}
\exo7qcmid{lille-2.integrales.029}

\contenu{
\begin{question}
\qclassification{127.99}
\qtags{motcle=changements de variables}
Parmi les affirmations suivantes, cocher celles qui sont vraies :
\begin{answers}
\bad{$\forall t\in \Rr\setminus\{-1\}$, $\displaystyle \frac{2t}{t^2+2t+1}=\frac{1}{t+1}-\frac{1}{(t+1)^2}$.}
\bad{$\displaystyle \ln (t+1)+\frac{1}{t+1}$ est une primitive de $\displaystyle \frac{2t}{t^2+2t+1}$ sur $]-1,+\infty[$.}
\good{Le changement de variable $t=\sqrt{x}$ donne $\displaystyle \int _0^1\frac{\mathrm{d}\, x}{x+2\sqrt{x}+1}=\int _0^1\frac{2t\mathrm{d}\, t}{t^2+2t+1}$.}
\good{$\displaystyle \int _0^1\frac{\mathrm{d}\, x}{x+2\sqrt{x}+1}=2\ln 2-1$.}
\end{answers}
\vskip3mm
\begin{explanations}
Avec $t=\sqrt{x}\Rightarrow x=t^2$, on a : $\mathrm{d}x=2t\, \mathrm{d}t$, $t(0)=0$, $t(1)=1$ et
$$\displaystyle \int _0^1\frac{\mathrm{d}\, x}{x+1+2\sqrt{x}}=\int _0^1\frac{2t\mathrm{d}\, t}{t^2+2t+1}=\left[2\ln (t+1)+\frac{2}{t+1}\right]_{0}^1=2\ln 2-1.$$
\end{explanations}
\end{question}
}
